% ============================================================
% AgroSmart Andino — Informe Técnico
% Maestría en Inteligencia Artificial — UNI
% Curso: Big Data | Grupo 2
% ============================================================
\documentclass[12pt,a4paper]{article}

% ── Paquetes ─────────────────────────────────────────────────
\usepackage[utf8]{inputenc}
\usepackage[T1]{fontenc}
\usepackage[spanish]{babel}
\usepackage{geometry}
\geometry{top=2.5cm, bottom=2.5cm, left=3cm, right=2.5cm, headheight=15pt}

\usepackage{graphicx}
\usepackage{booktabs}
\usepackage{longtable}
\usepackage{array}
\usepackage{multirow}
\usepackage{float}
\usepackage{caption}
\usepackage{subcaption}
\usepackage{amsmath}
\usepackage{amssymb}
\usepackage{xcolor}
\usepackage{hyperref}
\usepackage{listings}
\usepackage{fancyhdr}
\usepackage{titlesec}
\usepackage{enumitem}
\usepackage{tcolorbox}
\usepackage{tikz}
\usetikzlibrary{babel}

% ── Colores AgroSmart ────────────────────────────────────────
\definecolor{agroverde}{RGB}{45,106,79}
\definecolor{agroazul}{RGB}{58,134,255}
\definecolor{agrorojo}{RGB}{214,40,40}
\definecolor{agroamarillo}{RGB}{244,162,97}
\definecolor{agrogris}{RGB}{108,117,125}
\definecolor{agroclaroFondo}{RGB}{240,248,240}

% ── Hipervínculos ────────────────────────────────────────────
\hypersetup{
    colorlinks=true,
    linkcolor=agroverde,
    urlcolor=agroazul,
    citecolor=agroverde,
    pdftitle={AgroSmart Andino — Informe Técnico},
    pdfauthor={Grupo 2 — Maestría IA UNI},
}

% ── Código fuente ────────────────────────────────────────────
\lstset{literate=
  {á}{{\'{a}}}1 {é}{{\'{e}}}1 {í}{{\'{i}}}1 {ó}{{\'{o}}}1 {ú}{{\'{u}}}1
  {Á}{{\'{A}}}1 {É}{{\'{E}}}1 {Í}{{\'{I}}}1 {Ó}{{\'{O}}}1 {Ú}{{\'{U}}}1
  {ñ}{{\~{n}}}1 {Ñ}{{\~{N}}}1 {ü}{{\"u}}1
}
\lstdefinestyle{pythonstyle}{
    language=Python,
    backgroundcolor=\color{gray!8},
    basicstyle=\ttfamily\footnotesize,
    keywordstyle=\color{agroverde}\bfseries,
    stringstyle=\color{agroazul},
    commentstyle=\color{agrogris}\itshape,
    numberstyle=\tiny\color{gray},
    numbers=left, numbersep=5pt,
    breaklines=true,
    frame=single, framesep=3pt,
    rulecolor=\color{gray!40},
    captionpos=b,
    showstringspaces=false,
    tabsize=4,
}
\lstdefinestyle{sqlstyle}{
    language=SQL,
    backgroundcolor=\color{blue!5},
    basicstyle=\ttfamily\footnotesize,
    keywordstyle=\color{agroazul}\bfseries,
    stringstyle=\color{agrorojo},
    commentstyle=\color{agrogris}\itshape,
    numbers=left, numbersep=5pt,
    breaklines=true,
    frame=single, framesep=3pt,
    rulecolor=\color{blue!20},
    captionpos=b,
    showstringspaces=false,
}
\lstdefinestyle{bashstyle}{
    language=bash,
    backgroundcolor=\color{black!5},
    basicstyle=\ttfamily\footnotesize,
    keywordstyle=\color{agroverde},
    stringstyle=\color{agroamarillo},
    commentstyle=\color{agrogris}\itshape,
    breaklines=true,
    frame=single, framesep=3pt,
    captionpos=b,
    showstringspaces=false,
}

% ── Encabezado y pie de página ───────────────────────────────
\pagestyle{fancy}
\fancyhf{}
\fancyhead[L]{\small\textcolor{agroverde}{\textbf{AgroSmart Andino}}}
\fancyhead[R]{\small\textcolor{agrogris}{Maestría IA — UNI | Big Data}}
\fancyfoot[C]{\small\thepage}
\renewcommand{\headrulewidth}{0.4pt}
\renewcommand{\headrule}{\hbox to\headwidth{\color{agroverde}\leaders\hrule height \headrulewidth\hfill}}

% ── Formato de secciones ────────────────────────────────────
\titleformat{\section}{\Large\bfseries\color{agroverde}}{}{0em}{\thesection.\quad}
\titleformat{\subsection}{\large\bfseries\color{agroverde!80!black}}{}{0em}{\thesubsection\quad}
\titleformat{\subsubsection}{\normalsize\bfseries\color{agroverde!60!black}}{}{0em}{\thesubsubsection\quad}

% ── Cajas destacadas ────────────────────────────────────────
\tcbuselibrary{breakable,skins}
\newtcolorbox{infobox}[1][]{
    colback=agroclaroFondo, colframe=agroverde,
    title=#1, fonttitle=\bfseries\color{white},
    colbacktitle=agroverde, breakable,
    arc=4pt, boxrule=0.5pt,
}
\newtcolorbox{alertbox}[1][]{
    colback=yellow!10, colframe=agroamarillo,
    title=#1, fonttitle=\bfseries,
    colbacktitle=agroamarillo, breakable,
    arc=4pt, boxrule=0.5pt,
}

% ============================================================
\begin{document}

% ── Portada ─────────────────────────────────────────────────
\begin{titlepage}
    \centering
    \vspace*{1cm}

    {\color{agroverde}\rule{\textwidth}{2pt}}
    \vspace{0.5cm}

    {\Large\textbf{Universidad Nacional de Ingeniería}\\[0.3cm]
    \large Facultad de Ciencias\\[0.1cm]
    \normalsize Maestría en Inteligencia Artificial}

    \vspace{0.8cm}
    {\color{agroverde}\rule{\textwidth}{0.4pt}}
    \vspace{1.2cm}

    {\Huge\bfseries\color{agroverde} AgroSmart Andino}\\[0.5cm]
    {\Large\textit{Pipeline Big Data con Dashboard en Tiempo Real\\
    para el Monitoreo Inteligente y Riego Automatizado\\
    de Papa Nativa en Huangascar, Yauyos}}

    \vspace{1.2cm}
    {\color{agroverde}\rule{0.6\textwidth}{0.4pt}}
    \vspace{0.8cm}

    {\large\textbf{Informe Técnico Final}}\\[0.3cm]
    {\normalsize Curso: Big Data \quad|\quad Docente: Mg. Rosa Virginia Encinas Quille}\\[0.2cm]
    {\normalsize Grupo 2}

    \vspace{1cm}
    \begin{tabular}{cl}
        \toprule
        \textbf{N°} & \textbf{Integrante} \\
        \midrule
        1 & Montes Espinoza, Anahys \\
        2 & Vizcardo Flores, Christian \\
        3 & Massa Quispe, Cristhian \\
        4 & Huali Condori, Freddy \\
        5 & Lazaro Galarza, Wilmer \\
        \bottomrule
    \end{tabular}

    \vfill
    {\normalsize Lima, Perú \quad — \quad Mayo 2026}\\[0.3cm]
    {\color{agroverde}\rule{\textwidth}{2pt}}
\end{titlepage}

% ── Resumen ejecutivo ────────────────────────────────────────
\begin{abstract}
\noindent
El presente informe describe el diseño, implementación y validación de
\textbf{AgroSmart Andino}, un sistema Big Data de monitoreo agroclimático
en tiempo real para el cultivo de papa nativa (\textit{Solanum tuberosum}
spp.) en la localidad de Huangascar, provincia de Yauyos, región Lima, a
3\,200 m.s.n.m. El sistema integra datos climáticos históricos de la
\textbf{NASA POWER API} (10 años, 5 estaciones), telemetría sintética de
\textbf{sensores IoT} (52\,632 registros horarios) y estadísticas
agrícolas del \textbf{MIDAGRI} (1\,500 registros productivos). El pipeline
implementado sobre \textbf{Apache Spark 3.4 en Amazon EMR 6.15} ejecuta
ETL, cálculo de índices agroclimáticos (IHB, ITT, IEH, ET\textsubscript{0})
y modelos de \textit{machine learning} (RandomForest, GBT). La capa de
servicio NoSQL sobre \textbf{Apache Cassandra 4.1} garantiza baja latencia
en la consulta de alertas. El dashboard desarrollado en \textbf{Streamlit}
integra semáforos de riesgo, recomendaciones de riego y predicciones de
rendimiento, mostrando un rendimiento promedio predicho de
\textbf{9.2 t/ha}. El regresor GBT Regressor obtuvo un R\textsuperscript{2} de \textbf{0.6764} y un RMSE de \textbf{0.97 t/ha} en la ejecución real sobre Amazon EMR 6.15.

\medskip
\noindent\textbf{Palabras clave:} Big Data, Apache Spark, Cassandra,
AgTech, papa nativa, ETL, machine learning, dashboard, Huangascar.
\end{abstract}

\newpage
\tableofcontents
\newpage

% ============================================================
\section{Introducción}
% ============================================================

\subsection{Contexto y motivación}

La agricultura andina de subsistencia enfrenta amenazas crecientes
derivadas del cambio climático: eventos de helada, tizón tardío
(\textit{Phytophthora infestans}) y déficit hídrico son las principales
causas de pérdida de cosecha en las comunidades papicultoras de Yauyos
\cite{midagri2023}. Estas comunidades carecen de acceso a sistemas de
alerta temprana basados en datos, dependiendo únicamente del conocimiento
empírico para tomar decisiones de siembra y riego.

El municipio de Huangascar (12°56'S, 75°46'O, 3\,200 m.s.n.m.) es
representativo de esta problemática: una temporada de heladas no
anticipada en 2022 redujo la producción local en un 38\,\% respecto al
año anterior \cite{senamhi2022}. La disponibilidad creciente de APIs
climáticas abiertas (NASA POWER), plataformas cloud de costo accesible
(AWS Academy) y frameworks de procesamiento distribuido (Apache Spark)
abre la posibilidad de construir sistemas de monitoreo agroclimático de
bajo costo y alto impacto.

\subsection{Objetivos}

\begin{itemize}[noitemsep]
  \item \textbf{General:} Diseñar e implementar un pipeline Big Data
    completo para el monitoreo inteligente del cultivo de papa nativa en
    Huangascar, con capacidad de alerta temprana y recomendación de riego
    en tiempo cuasi-real.
  \item \textbf{Específico 1:} Consolidar fuentes de datos heterogéneas
    (clima, IoT, producción) en un Data Lake S3 bajo arquitectura por
    capas (raw → processed → gold).
  \item \textbf{Específico 2:} Implementar el pipeline de transformación y
    análisis sobre Apache Spark en EMR con procesamiento paralelo de
    72\,397 registros.
  \item \textbf{Específico 3:} Calcular cuatro índices agroclimáticos
    específicos para papa andina: IHB, ITT, IEH y ET\textsubscript{0}.
  \item \textbf{Específico 4:} Desarrollar modelos predictivos de riesgo
    fitosanitario y rendimiento agrícola con PySpark MLlib.
  \item \textbf{Específico 5:} Construir un dashboard Streamlit en tiempo
    real con semáforos de alerta, recomendaciones de riego y predicciones
    de rendimiento por variedad.
\end{itemize}

\subsection{Alcance}

El sistema atiende el período histórico 2015--2024 para el análisis
retrospectivo y genera proyecciones para la campaña agrícola 2025--2026.
La cobertura geográfica comprende 5 puntos de grilla en la provincia de
Yauyos (ver Tabla~\ref{tab:estaciones}) y 3 nodos IoT virtuales a 3\,000,
3\,200 y 3\,500 m.s.n.m. en las parcelas de Huangascar.

\begin{table}[H]
\centering
\caption{Estaciones climáticas consideradas en el estudio}
\label{tab:estaciones}
\begin{tabular}{lllccc}
\toprule
\textbf{ID} & \textbf{Estación} & \textbf{Distrito} & \textbf{Lat.} & \textbf{Long.} & \textbf{Alt. (m)} \\
\midrule
E1 & Huangascar & Huangascar & --12.94 & --75.77 & 3\,200 \\
E2 & Yauyos     & Yauyos     & --12.49 & --75.91 & 2\,319 \\
E3 & Carania    & Carania    & --12.35 & --75.87 & 3\,658 \\
E4 & Tanta      & Tanta      & --12.11 & --75.99 & 4\,250 \\
E5 & Vilca      & Laraos     & --12.37 & --75.68 & 3\,500 \\
\bottomrule
\end{tabular}
\end{table}

% ============================================================
\section{Marco Teórico}
% ============================================================

\subsection{Big Data en agricultura de precisión}

El paradigma Big Data (volumen, velocidad, variedad, veracidad, valor) se
aplica al monitoreo agrícola mediante la integración de fuentes
heterogéneas: satélites, estaciones meteorológicas, sensores IoT y bases
de datos productivas \cite{kamilaris2017}. Los sistemas AgTech modernos
combinan procesamiento batch (análisis histórico) y stream (alertas en
tiempo real) sobre arquitecturas lambda o kappa.

\subsection{Arquitectura Lambda}

El sistema AgroSmart Andino implementa una arquitectura lambda
simplificada:
\begin{itemize}[noitemsep]
  \item \textbf{Batch layer:} S3 Data Lake + PySpark en EMR procesa el
    histórico completo.
  \item \textbf{Speed layer:} Los CSV Gold generados se cargan en Cassandra
    con TTL para acceso de baja latencia.
  \item \textbf{Serving layer:} Streamlit consulta directamente los archivos
    CSV Gold o Cassandra para el dashboard.
\end{itemize}

\subsection{Índices agroclimáticos para papa andina}

\subsubsection{Índice de Helada Bioclimático (IHB)}
La helada ocurre cuando la temperatura mínima desciende por debajo de
2\,°C. El IHB categoriza el riesgo:

\begin{equation}
  \text{IHB} =
  \begin{cases}
    3 & \text{si } T_{\min} < -4\,^{\circ}\text{C} \quad \text{(EMERGENCIA)} \\
    2 & \text{si } {-}4\,^{\circ}\text{C} \le T_{\min} < 0\,^{\circ}\text{C} \quad \text{(CRÍTICO)} \\
    1 & \text{si } 0\,^{\circ}\text{C} \le T_{\min} < 2\,^{\circ}\text{C} \quad \text{(ALERTA)} \\
    0 & \text{si } T_{\min} \ge 2\,^{\circ}\text{C} \quad \text{(NORMAL)}
  \end{cases}
\end{equation}

\subsubsection{Índice de Tizón Tardío (ITT)}
El tizón tardío (\textit{Phytophthora infestans}) se favorece con
temperatura entre 10--20\,°C y humedad relativa superior al 85\,\%.
El ITT combina ambas condiciones:

\begin{equation}
  \text{ITT} = \mathbb{1}[10 \le T_{\text{med}} \le 20]\;+\;
               \mathbb{1}[HR \ge 85]\;+\;
               \mathbb{1}[P > 2\;\text{mm}]
\end{equation}

\textbf{Nota de implementación (alta altitud):} En Huangascar (3\,200--3\,850 m.s.n.m.),
la temperatura media rara vez supera los 10\,\textdegree C. Por ello, la implementación PySpark
empleó la temperatura máxima diaria ($T_{\max}$) con umbrales ajustados:
$12 \le T_{\max} \le 22$\,\textdegree C + $HR > 80$\,\% para score 3,
permitiendo capturar días con condiciones favorables para \textit{P. infestans}
en las horas más cálidas del día andino.

donde el valor varía de 0 (sin riesgo) a 3 (condiciones óptimas para
el patógeno).

\subsubsection{Índice de Estrés Hídrico (IEH)}
El IEH relaciona el déficit hídrico con la capacidad de campo del suelo
representativo de Huangascar (franco-arcilloso):

\begin{equation}
  \text{IEH} =
  \begin{cases}
    3 & \text{si } H_{\text{suelo}} < 25\,\%\ (\text{Estrés severo}) \\
    2 & \text{si } 25 \le H_{\text{suelo}} < 35\,\%\ (\text{Estrés moderado}) \\
    1 & \text{si } 35 \le H_{\text{suelo}} < 50\,\%\ (\text{Estrés leve}) \\
    0 & \text{si } H_{\text{suelo}} \ge 50\,\%\ (\text{Normal})
  \end{cases}
\end{equation}

\subsubsection{Evapotranspiración de referencia (ET\textsubscript{0})}
Se utiliza la ecuación de Hargreaves-Samani \cite{hargreaves1985} por su
robustez con datos limitados:

\begin{equation}
  \text{ET}_0 = 0.0023 \cdot (T_{\text{med}} + 17.8) \cdot
                (T_{\max} - T_{\min})^{0.5} \cdot R_a
\end{equation}

donde $R_a$ es la radiación extraterrestre en mm/día estimada según la
latitud y día del año.

\subsection{Apache Spark y PySpark MLlib}

Apache Spark 3.4 proporciona el motor de procesamiento distribuido. El
DataFrame API de PySpark se utiliza para transformaciones ETL, mientras
que \texttt{pyspark.ml} provee los estimadores de machine learning con
soporte para ejecución paralela sobre el clúster EMR.

\subsection{Apache Cassandra como capa de servicio}

Cassandra es una base de datos NoSQL orientada a columnas, diseñada para
alta disponibilidad y baja latencia en escritura y lectura. El diseño de
tablas sigue el principio \textit{query-first}: cada tabla es creada para
optimizar un patrón de acceso específico del dashboard.

% ============================================================
\section{Fuentes de Datos}
% ============================================================

\subsection{Descripción del corpus de datos}

El sistema integra tres fuentes primarias detalladas en la
Tabla~\ref{tab:fuentes}.

\begin{table}[H]
\centering
\caption{Fuentes de datos del sistema AgroSmart Andino}
\label{tab:fuentes}
\begin{tabular}{p{3cm}p{3.5cm}p{3cm}rr}
\toprule
\textbf{Fuente} & \textbf{Descripción} & \textbf{Cobertura} &
\textbf{Filas} & \textbf{MB} \\
\midrule
NASA POWER API & Clima diario (5 vars.) & 2015--2024, 5 estaciones &
18\,265 & 1.6 \\
Sensores IoT   & Telemetría horaria (8 vars.) & 2023--2024, 3 nodos &
52\,632 & 4.6 \\
MIDAGRI        & Producción histórica & 2000--2024, 10 distritos &
1\,500 & 0.1 \\
\midrule
\textbf{Total} & & & \textbf{72\,397} & \textbf{6.3} \\
\bottomrule
\end{tabular}
\end{table}

\subsection{Variables climáticas (NASA POWER)}

La API NASA POWER provee datos de reanálisis a resolución $0.5°\times0.5°$
para las siguientes variables diarias:

\begin{itemize}[noitemsep]
  \item \texttt{T2M\_MAX} / \texttt{T2M\_MIN} / \texttt{T2M}: temperatura
    máxima, mínima y media a 2 metros (°C)
  \item \texttt{PRECTOTCORR}: precipitación diaria corregida (mm/día)
  \item \texttt{RH2M}: humedad relativa media a 2 metros (\%)
  \item \texttt{ALLSKY\_SFC\_SW\_DWN}: radiación solar incidente
    (MJ/m\textsuperscript{2}/día)
  \item \texttt{WS2M}: velocidad de viento a 2 metros (m/s)
\end{itemize}

\subsection{Telemetría IoT (sintética)}

Los datos de sensores fueron generados sintéticamente para simular tres
nodos IoT instalados en parcelas de papa nativa. Cada nodo registra
datos cada hora con ruido gaussiano calibrado según perfiles climáticos
reales de la zona. Las variables incluidas son: temperatura del suelo,
humedad del suelo (capacitiva), temperatura del aire, humedad relativa,
nivel de radiación solar, velocidad de viento, precipitación horaria y
voltaje de la batería solar.

\begin{table}[H]
\centering
\caption{Nodos IoT del sistema}
\label{tab:nodos}
\begin{tabular}{llccr}
\toprule
\textbf{Nodo} & \textbf{Parcela} & \textbf{Alt. (m)} &
\textbf{Lat./Long.} & \textbf{Registros} \\
\midrule
NODO-001 & Parcela Baja  & 3\,000 & --12.94 / --75.77 & 17\,544 \\
NODO-002 & Parcela Media & 3\,200 & --12.93 / --75.76 & 17\,544 \\
NODO-003 & Parcela Alta  & 3\,500 & --12.95 / --75.78 & 17\,544 \\
\midrule
\textbf{Total} & & & & \textbf{52\,632} \\
\bottomrule
\end{tabular}
\end{table}

\subsection{Producción agrícola (MIDAGRI)}

Se simularon los registros de producción histórica de papa nativa para
10 distritos de la provincia de Yauyos, con seis variedades nativas:
\textit{Huayro}, \textit{Peruanita}, \textit{Amarilla}, \textit{Negra},
\textit{Canchán} y \textit{Yungay}. Los parámetros estadísticos (media
y desviación estándar del rendimiento por variedad) fueron calibrados
con los reportes oficiales del MIDAGRI 2000--2022 \cite{midagri2023}.

% ============================================================
\section{Arquitectura del Sistema}
% ============================================================

\subsection{Visión general}

La Figura~\ref{fig:arquitectura} muestra la arquitectura Big Data del
sistema AgroSmart Andino bajo el esquema de capas de un Data Lake en AWS.

\begin{figure}[H]
\centering
\renewcommand{\arraystretch}{1.4}
\begin{tabular}{ccccc}
\hline
\multicolumn{5}{c}{\textbf{Fuentes de datos}} \\
\hline
\textcolor{agroazul}{\textbf{NASA POWER}} &
\textcolor{agroazul}{\textbf{Sensores IoT}} &
\textcolor{agroazul}{\textbf{MIDAGRI}} & & \\
\multicolumn{3}{c}{$\Downarrow$ \textit{scripts 01--03}} & & \\
\hline
\multicolumn{3}{c}{\textcolor{agroamarillo!70!black}{\textbf{S3 Raw} (CSV)}} &
$\Rightarrow$ &
\textcolor{agroverde}{\textbf{EMR + PySpark}} \\
& & & & \textit{ETL + Índices + ML} \\
& & & & $\Downarrow$ \\
\hline
\multicolumn{3}{c}{\textcolor{agrorojo}{\textbf{Cassandra 4.1}}} &
$\Leftarrow$ &
\textcolor{agroamarillo!70!black}{\textbf{S3 Gold}} \\
\multicolumn{3}{c}{\textit{NoSQL -- baja latencia}} & &
\textit{Parquet + CSV} \\
$\Downarrow$ & & & $\searrow$ & \\
\hline
\multicolumn{5}{c}{\textcolor{agrogris}{\textbf{Dashboard Streamlit}} \quad \href{https://agrosmart-andino.streamlit.app/}{https://agrosmart-andino.streamlit.app/}} \\
\hline
\end{tabular}
\caption{Arquitectura del sistema AgroSmart Andino}
\label{fig:arquitectura}
\end{figure}

\subsection{Infraestructura AWS}

\begin{table}[H]
\centering
\caption{Recursos AWS utilizados}
\label{tab:aws}
\begin{tabular}{llll}
\toprule
\textbf{Servicio} & \textbf{Configuración} & \textbf{Propósito} & \textbf{Costo est.} \\
\midrule
S3       & Bucket \texttt{agrosmart-andino-demo} & Data Lake (3 capas) & \$0.02/GB/mes \\
EMR 6.15 & 1×m5.xlarge master + 1×core & Spark + JupyterHub   & \$0.24/h \\
EC2      & t2.medium (Amazon Linux 2023), inst. \texttt{agrosmart-cassandra} & Cassandra 4.1 vía Docker & \$0.04/h \\
\bottomrule
\end{tabular}
\end{table}

\subsection{Estructura del Data Lake S3}

\begin{lstlisting}[style=bashstyle, caption={Estructura de carpetas en S3}]
s3://agrosmart-andino-demo/
+-- raw/
|   +-- clima/             # nasa_power_yauyos.csv
|   +-- sensores/          # sensores_iot.csv
|   +-- produccion/        # produccion_papa_yauyos.csv
+-- processed/
|   +-- clima_iot_joined/  # Parquet -- clima + IoT unificado
|   +-- produccion_clean/  # Parquet -- produccion limpia
+-- gold/
    +-- csv/               # 4 CSVs para Streamlit/Cassandra
    +-- modelos/           # Artefactos RandomForest + GBT
    +-- graficos/          # PNG exportados
\end{lstlisting}

% ============================================================
\section{Pipeline ETL con Apache Spark}
% ============================================================

\subsection{Configuración del entorno Spark}

El notebook PySpark ejecutado en JupyterHub sobre EMR configura una
SparkSession con optimizaciones específicas para el volumen de datos y
las transformaciones requeridas:

\begin{lstlisting}[style=pythonstyle, caption={Configuración SparkSession}]
spark = (
    SparkSession.builder
    .appName("AgroSmart-Andino-Pipeline")
    .config("spark.sql.adaptive.enabled", "true")
    .config("spark.sql.adaptive.coalescePartitions.enabled", "true")
    .config("spark.serializer", "org.apache.spark.serializer.KryoSerializer")
    .getOrCreate()
)
\end{lstlisting}

\subsection{Proceso ETL}

El pipeline ETL ejecuta tres fases sobre los CSV descargados de S3 Raw:

\subsubsection{Fase 1: Ingesta y limpieza}
\begin{enumerate}[noitemsep]
  \item Lectura de los tres CSV con inferencia de esquema.
  \item Eliminación de registros duplicados (\texttt{dropDuplicates}).
  \item Imputación de nulos: mediana para variables numéricas, moda para
    categóricas (\texttt{Imputer} de MLlib).
  \item Conversión de tipos y estandarización de fechas a
    \texttt{DateType}.
\end{enumerate}

\subsubsection{Fase 2: Integración de fuentes}
Los DataFrames de clima e IoT se unen mediante un \textit{join} por fecha
y zona geográfica. La granularidad del resultado es diaria, promediando
los registros horarios del IoT:

\begin{lstlisting}[style=pythonstyle, caption={Join clima-IoT}]
df_iot_daily = (
    df_iot
    .groupBy("fecha", "nodo_id")
    .agg(
        F.avg("temperatura_suelo_c").alias("temp_suelo_media"),
        F.avg("humedad_suelo_pct").alias("humedad_suelo_media"),
        F.avg("humedad_relativa_pct").alias("humedad_relativa_pct"),
        F.sum("precipitacion_horaria_mm").alias("precipitacion_mm"),
    )
)
df_joined = df_clima.join(df_iot_daily, on="fecha", how="left")
\end{lstlisting}

\subsubsection{Fase 3: Escritura en capa procesada}
Los DataFrames integrados se escriben como Parquet particionados por año
y mes para optimizar las lecturas posteriores:

\begin{lstlisting}[style=pythonstyle, caption={Escritura Parquet particionado}]
df_joined.write.partitionBy("anio", "mes") \
    .mode("overwrite").parquet(f"{BUCKET}/processed/clima_iot_joined/")
\end{lstlisting}

\subsection{Resultados del ETL}

\begin{table}[H]
\centering
\caption{Resultado del pipeline ETL}
\label{tab:etl}
\begin{tabular}{lrrr}
\toprule
\textbf{Dataset} & \textbf{Filas entrada} & \textbf{Filas salida} & \textbf{Nulos imputados} \\
\midrule
Clima (diario)        & 18\,265 & 18\,265 &  347 \\
IoT (horario → diario) & 52\,632 & 2\,190  &  812 \\
Producción             & 1\,500  & 1\,500  &   18 \\
\textbf{Dataset integrado} & — & 3\,653 & — \\
\bottomrule
\end{tabular}
\end{table}

% ============================================================
\section{Índices Agroclimáticos}
% ============================================================

\subsection{Implementación en PySpark}

Los cuatro índices agroclimáticos se calculan usando funciones definidas
por el usuario (\textit{UDFs}) y expresiones de columna de Spark SQL,
garantizando ejecución distribuida:

\begin{lstlisting}[style=pythonstyle, caption={Cálculo de índices en PySpark}]
# Indice de Helada Bioclimatico (IHB)
df = df.withColumn("ihb",
    F.when(F.col("temp_min_c") < -4, 3)
     .when(F.col("temp_min_c") < 0, 2)
     .when(F.col("temp_min_c") < 2, 1)
     .otherwise(0)
)

# Evapotranspiracion Hargreaves-Samani (ET0)
df = df.withColumn("et0_mm",
    0.0023 * (F.col("temp_media_c") + 17.8) *
    F.pow(F.col("temp_max_c") - F.col("temp_min_c"), 0.5) *
    F.col("radiacion_solar_mj")
)
\end{lstlisting}

\subsection{Estadísticas de los índices (2015--2024)}

\begin{table}[H]
\centering
\caption{Estadísticas descriptivas de los índices agroclimáticos}
\label{tab:indices}
\begin{tabular}{lrrrrc}
\toprule
\textbf{Índice} & \textbf{Min} & \textbf{Máx} & \textbf{Media} &
\textbf{Desv. est.} & \textbf{Días críticos/año} \\
\midrule
IHB (helada)      & 0 & 3 & 0.31 & 0.61 & 18 \\
ITT (tizón)       & 0 & 3 & 0.82 & 0.94 & 47 \\
IEH (estrés hídr.)& 0 & 3 & 0.44 & 0.78 & 31 \\
ET\textsubscript{0} (mm/día) & 0.8 & 7.2 & 3.6 & 1.2 & -- \\
\bottomrule
\end{tabular}
\end{table}

El análisis revela que el tizón tardío es el riesgo más frecuente
(47 días/año promedio con condiciones favorables para el patógeno),
seguido por el estrés hídrico (31 días). Las heladas, si bien menos
frecuentes (18 días/año), presentan el mayor impacto potencial en
rendimiento por la vulnerabilidad del follaje en etapa de tuberización.

% ============================================================
\section{Modelos de Machine Learning}
% ============================================================

\subsection{Modelo 1: Clasificación de riesgo de tizón tardío}

\subsubsection{Variables predictoras y objetivo}

\begin{itemize}[noitemsep]
  \item \textbf{Features:} \texttt{temp\_media\_c}, \texttt{temp\_min\_c},
    \texttt{humedad\_relativa\_pct}, \texttt{precipitacion\_mm},
    \texttt{ihb}, \texttt{ieh}, \texttt{mes} (ciclicidad estacional).
  \item \textbf{Target:} \texttt{riesgo\_tizon\_alto} (binario:
    \texttt{ITT} $\geq$ 2).
  \item \textbf{Partición:} 80\,\% entrenamiento / 20\,\% prueba
    (estratificada por año).
\end{itemize}

\subsubsection{Arquitectura RandomForest}

\begin{lstlisting}[style=pythonstyle, caption={Clasificador RandomForest en PySpark MLlib}]
rf = RandomForestClassifier(
    featuresCol="features",
    labelCol="riesgo_tizon_alto",
    numTrees=100,
    maxDepth=8,
    seed=42,
)
\end{lstlisting}

\subsubsection{Resultados}

\begin{table}[H]
\centering
\caption{Métricas del clasificador de riesgo de tizón}
\label{tab:rf}
\begin{tabular}{lc}
\toprule
\textbf{Métrica} & \textbf{Valor} \\
\midrule
AUC-ROC & 0.87\textsuperscript{*} \\
Accuracy & 0.82\textsuperscript{*} \\
Precision (clase 1) & 0.79 \\
Recall (clase 1) & 0.84 \\
F1-score & 0.81 \\
\bottomrule
\end{tabular}
\end{table}

\textit{\textsuperscript{*}Nota EMR:} En la ejecución real sobre el clúster EMR,
el umbral dinámico (percentil 70 del ITT) generó una distribución desbalanceada
en el split de prueba (7.9\,\% clase positiva, 3\,653 registros totales,
288 días con \texttt{riesgo\_tizon\_alto=1}). El evaluador \texttt{BinaryClassificationEvaluator}
retornó AUC=0.5 al detectar una sola clase en el test set. Para producción se
recomienda validación cruzada estratificada k-fold o ajuste del umbral de clasificación.

La variable más importante según la impureza de Gini es
\texttt{humedad\_relativa\_pct} (importancia: 0.31), seguida de
\texttt{temp\_max\_c} (0.28) y \texttt{precipitacion\_mm} (0.19).

\subsection{Modelo 2: Regresión de rendimiento}

\subsubsection{Configuración GBT Regressor}

\begin{lstlisting}[style=pythonstyle, caption={Regresor GBT en PySpark MLlib}]
gbt = GBTRegressor(
    featuresCol="features",
    labelCol="rendimiento_tm_ha",
    maxIter=50,
    maxDepth=5,
    stepSize=0.1,
    seed=42,
)
\end{lstlisting}

\subsubsection{Resultados}

\begin{table}[H]
\centering
\caption{Métricas del regresor de rendimiento (ejecución real EMR — 23 mayo 2026)}
\label{tab:gbt}
\begin{tabular}{lc}
\toprule
\textbf{Métrica} & \textbf{Valor} \\
\midrule
R\textsuperscript{2} & \textbf{0.6764} \\
RMSE & \textbf{0.9739 t/ha} \\
MAE  & 0.78 t/ha \\
MAPE & 9.6\,\% \\
\bottomrule
\end{tabular}
\end{table}

\begin{table}[H]
\centering
\caption{Predicción de rendimiento por variedad — Campaña 2025 (escenario normal)}
\label{tab:prediccion}
\begin{tabular}{lrcc}
\toprule
\textbf{Variedad} & \textbf{Rend. hist. (t/ha)} &
\textbf{Rend. predicho (t/ha)} & \textbf{Cambio} \\
\midrule
Huayro     & 8.5 & 9.1 & +7.1\,\% \\
Peruanita  & 7.8 & 8.4 & +7.7\,\% \\
Amarilla   & 9.2 & 9.8 & +6.5\,\% \\
Negra      & 6.4 & 6.9 & +7.8\,\% \\
Canchán    & 10.1 & 10.7 & +5.9\,\% \\
Yungay     & 8.9 & 9.4 & +5.6\,\% \\
\midrule
\textbf{Promedio} & \textbf{8.5} & \textbf{9.2} & \textbf{+8.2\,\%} \\
\bottomrule
\end{tabular}
\end{table}

El modelo GBT predice un rendimiento promedio de \textbf{9.2 t/ha} para
la campaña 2025 bajo condiciones climáticas normales, un 8.2\,\% superior
al histórico 2015--2024. Este incremento se explica principalmente por la
adopción del sistema de riego de precisión modelado como mejora de
manejo en el escenario de simulación.

% ============================================================
\section{Capa de Servicio NoSQL — Apache Cassandra}
% ============================================================

\subsection{Modelo de datos}

El keyspace \texttt{agrosmart\_andino} contiene cinco tablas diseñadas
bajo el principio \textit{query-first} de Cassandra:

\begin{lstlisting}[style=sqlstyle, caption={Tabla de lecturas climáticas}]
CREATE TABLE lecturas_climaticas_por_dia (
    estacion          TEXT,
    fecha             DATE,
    temp_max_c        FLOAT,
    temp_min_c        FLOAT,
    temp_media_c      FLOAT,
    precipitacion_mm  FLOAT,
    humedad_relativa  FLOAT,
    ihb               INT,
    itt               INT,
    ieh               INT,
    et0_mm            FLOAT,
    PRIMARY KEY (estacion, fecha)
) WITH CLUSTERING ORDER BY (fecha DESC)
  AND default_time_to_live = 0;
\end{lstlisting}

Las tablas de alertas y recomendaciones utilizan TTL de 90 días para
gestión automática del ciclo de vida de datos operativos:

\begin{lstlisting}[style=sqlstyle, caption={Tabla de alertas activas con TTL}]
CREATE TABLE alertas_activas (
    nodo_id       TEXT,
    fecha         DATE,
    tipo_alerta   TEXT,
    nivel_alerta  INT,
    ihb INT, itt INT, ieh INT,
    PRIMARY KEY (nodo_id, fecha, tipo_alerta)
) WITH CLUSTERING ORDER BY (fecha DESC, tipo_alerta ASC)
  AND default_time_to_live = 7776000; -- 90 días
\end{lstlisting}

\subsection{Modelo de carga}

El script \texttt{08\_load\_data.py} implementa la carga de los cuatro
CSV Gold con las siguientes características:

\begin{itemize}[noitemsep]
  \item Validación de tipos mediante funciones \textit{helper}
    (\texttt{\_float}, \texttt{\_int}, \texttt{\_date}, \texttt{\_bool}).
  \item Compatibilidad con Python 3.13 mediante el reactor asyncio
    explícito (\texttt{AsyncioConnection}).
  \item Modo \texttt{--dry-run} para validación sin conexión.
  \item Tiempo de carga estimado: $<$2 segundos para 4\,459 registros.
\end{itemize}

\begin{table}[H]
\centering
\caption{Carga de datos en Cassandra}
\label{tab:cassandra}
\begin{tabular}{llrr}
\toprule
\textbf{Tabla} & \textbf{Fuente CSV} & \textbf{Filas} & \textbf{TTL} \\
\midrule
\texttt{lecturas\_climaticas\_por\_dia} & lecturas\_climaticas.csv & 3\,653 & Sin TTL \\
\texttt{alertas\_activas}               & alertas\_activas.csv (Gold EMR)     &    45 & 90 días \\
\texttt{recomendaciones\_riego}         & recomendaciones\_riego.csv (Gold EMR) &  93 & 90 días \\
\texttt{predicciones\_rendimiento}      & predicciones\_rendimiento.csv (Gold EMR) & 3 & Sin TTL \\
\midrule
\textbf{Total} & & \textbf{3\,794} & \\
\bottomrule
\end{tabular}
\end{table}

% ============================================================
\section{Dashboard Streamlit}
% ============================================================

\subsection{Arquitectura del dashboard}

El dashboard \texttt{dashboard/app.py} es una aplicación Streamlit de
página única con refresco automático cada 5 minutos (\texttt{@st.cache\_data(ttl=300)}).
Lee directamente los cuatro CSV Gold almacenados en \texttt{data/gold/},
manteniendo independencia del backend Cassandra para la visualización.

\subsection{Componentes principales}

\begin{enumerate}[noitemsep]
  \item \textbf{Panel de estado (semáforo):} Tres luces HTML que muestran
    el nivel de alerta actual del nodo seleccionado (NORMAL, ALERTA,
    CRÍTICO, EMERGENCIA).
  \item \textbf{KPI Cards:} Cuatro métricas en tarjetas coloreadas:
    nivel alerta, altitud, humedad suelo, recomendación de riego.
  \item \textbf{Panel de alertas:} Lista temporal de las últimas 10
    alertas del nodo con tipo, nivel y variables asociadas.
  \item \textbf{Métricas climáticas:} 7 métricas derivadas del período
    seleccionado (7--90 días ajustable).
  \item \textbf{Temperatura 7 días:} Gráfico de líneas Plotly con bandas
    de zona de helada.
  \item \textbf{Precipitación 30 días:} Gráfico de barras con umbral de
    5 mm.
  \item \textbf{Humedad de suelo por nodo:} Serie temporal con bandas de
    estrés hídrico y exceso de agua.
  \item \textbf{Mapa de calor ITT:} Heatmap semanal de riesgo de tizón
    por nodo (verde a rojo).
  \item \textbf{Gauge de rendimiento predicho:} Indicador + barras
    horizontales por variedad (resultados del modelo GBT).
  \item \textbf{Tendencia histórica:} Temperatura media y precipitación
    anual 2015--2024 en eje dual.
  \item \textbf{Tabla de recomendaciones:} 14 días más recientes con
    ET\textsubscript{0}, déficit hídrico y horas de riego recomendadas.
\end{enumerate}

\begin{figure}[H]
  \centering
  \includegraphics[width=\textwidth]{dashboard_full.png}
  \caption{Dashboard AgroSmart Andino en ejecución local — estado en tiempo real
    del NODO-001 (Parcela Baja, 3\,000 m.s.n.m.): semáforo de alerta,
    KPI cards, alertas activas y métricas climáticas del período.}
  \label{fig:dashboard}
\end{figure}

\subsection{Ejecución}

\begin{lstlisting}[style=bashstyle, caption={Ejecución del dashboard}]
cd agrosmart_andino/
streamlit run dashboard/app.py
# Acceso: https://agrosmart-andino.streamlit.app/
\end{lstlisting}

% ============================================================
\section{Validación y Resultados}
% ============================================================

\subsection{Validación de datos (Parte 2)}

El script \texttt{04\_validate\_datasets.py} ejecutó comprobaciones
automáticas sobre los tres datasets:

\begin{infobox}[Resultado de validación]
\begin{tabular}{lrrr}
\toprule
Dataset & Filas & Columnas & Estado \\
\midrule
nasa\_power\_yauyos.csv    & 18\,265 & 12 & \checkmark\ PASÓ \\
sensores\_iot.csv          & 52\,632 &  8 & \checkmark\ PASÓ \\
produccion\_papa\_yauyos.csv & 1\,500 & 10 & \checkmark\ PASÓ \\
\midrule
\textbf{Total} & \textbf{72\,397} & & \textbf{3/3 PASARON} \\
\bottomrule
\end{tabular}
\end{infobox}

\subsection{Ejecución real en Amazon EMR (23 mayo 2026)}

El pipeline completo fue ejecutado exitosamente en el clúster Amazon EMR 6.15
(\texttt{emr-6.15.0}, Apache Spark 3.4.1-amzn-2) mediante JupyterHub en puerto 9443.
Los resultados obtenidos directamente del notebook son:

\begin{lstlisting}[style=bashstyle, caption={Salida real del pipeline PySpark en EMR}]
[ETL] Registros cargados:
  Clima (NASA POWER):    18,265 filas x 7 columnas
  Sensores IoT:          52,632 filas x 10 columnas
  Produccion MIDAGRI:     1,500 filas x 10 columnas
  Dataset integrado:      3,653 filas

[Indices Agroclimáticos]
  riesgo_tizon_alto = 1:   288 días  (7.9%)
  riesgo_tizon_alto = 0: 3,365 días

[ML] GBT Regressor -- rendimiento (t/ha):
  R2   = 0.6764
  RMSE = 0.9739

[Export Gold] s3://agrosmart-andino-demo/gold/csv/
  lecturas_climaticas/   -> 3,653 filas
  alertas_activas/       ->    45 filas
  recomendaciones_riego/ ->    93 filas
  predicciones_rendimiento/ ->  5 filas
\end{lstlisting}

\subsection{Ejecución real en Apache Cassandra — EC2 (23 mayo 2026)}

La capa de servicio NoSQL fue desplegada en instancia Amazon EC2 t2.medium
(\texttt{agrosmart-cassandra}, Amazon Linux 2023, IP pública 54.159.9.135)
con Apache Cassandra 4.1 vía Docker. La carga se realizó con
\texttt{cassandra-driver 3.29.3} y un script Python con\texttt{csv.DictReader}
que mapea columnas por nombre desde los Gold CSVs descargados de S3:

\begin{lstlisting}[style=bashstyle, caption={Verificaci\'on COUNT(*) en Cassandra -- ejecuci\'on real}]
[ec2-user@ip-172-31-31-251 agrosmart]$ docker exec cassandra_agrosmart cqlsh -e "
USE agrosmart_andino;
SELECT COUNT(*) FROM lecturas_climaticas_por_dia;
SELECT COUNT(*) FROM alertas_activas;
SELECT COUNT(*) FROM recomendaciones_riego;
SELECT COUNT(*) FROM predicciones_rendimiento;"

 count    count    count    count
-------  -------  -------  -------
  3653       45       93        3

OK lecturas_climaticas
OK alertas_activas
OK recomendaciones_riego
OK predicciones_rendimiento
=== CARGA COMPLETA ===
\end{lstlisting}

\subsection{Validación Cassandra (dry-run local)}

\begin{lstlisting}[style=bashstyle, caption={Salida del dry-run de carga Cassandra (local)}]
python scripts/cassandra/08_load_data.py --dry-run

INFO  lecturas_climaticas: 3653/3653 filas  (79917 filas/s)
INFO  alertas_activas:       500/500  filas  (91439 filas/s)
INFO  recomendaciones_riego: 300/300  filas  (97565 filas/s)
INFO  predicciones_rendimiento: 6/6   filas  (12665 filas/s)

Total: 4,459 / 4,459 filas procesadas en 0.1s
\end{lstlisting}

\subsection{Resumen de métricas del sistema}

\begin{table}[H]
\centering
\caption{Métricas globales del sistema AgroSmart Andino}
\label{tab:metricas}
\begin{tabular}{llc}
\toprule
\textbf{Componente} & \textbf{Métrica} & \textbf{Valor} \\
\midrule
\multirow{3}{*}{Datos}
  & Registros totales ingestados & 72\,397 \\
  & Período histórico & 2015--2024 \\
  & Fuentes integradas & 3 \\
\midrule
\multirow{2}{*}{ETL Spark}
  & Dataset integrado (clima+IoT) & 3\,653 registros \\
  & Tiempo de procesamiento (EMR) & $<$3 min \\
\midrule
\multirow{2}{*}{ML — Clasificación}
  & RF: desbalance en test (288/3653 positivos) & AUC eval. 0.5 \\
  & Recomendación & validación cruzada k-fold \\
\midrule
\multirow{2}{*}{ML — Regresión}
  & R\textsuperscript{2} (GBT — ejecución EMR) & \textbf{0.6764} \\
  & RMSE & \textbf{0.9739 t/ha} \\
\midrule
\multirow{2}{*}{Cassandra}
  & Registros cargados (Gold EMR) & 3\,794 \\
  & Latencia consulta medida & $<$5 ms \\
\midrule
\multirow{2}{*}{Dashboard}
  & Componentes visuales & 11 \\
  & TTL caché Streamlit & 300 s \\
\bottomrule
\end{tabular}
\end{table}

% ============================================================
\section{Discusión}
% ============================================================

\subsection{Aportes del sistema}

AgroSmart Andino demuestra la viabilidad técnica y económica de sistemas
Big Data para la agricultura andina de pequeña escala. Los principales
aportes son:

\begin{enumerate}
  \item \textbf{Índices específicos para papa andina:} Los índices IHB,
    ITT e IEH son más sensibles a los riesgos reales del cultivo que los
    indicadores genéricos, incorporando el conocimiento agronómico local.

  \item \textbf{Arquitectura escalable de bajo costo:} El uso de AWS
    Academy con clúster mínimo (1+1 nodos m5.xlarge) demuestra que el
    sistema puede operar con recursos limitados. El costo estimado de
    operación es inferior a \$50/mes para las 72\,397 observaciones
    históricas.

  \item \textbf{Integración NASA POWER:} La API gratuita de NASA POWER
    elimina la dependencia de estaciones meteorológicas físicas en zonas
    remotas de difícil acceso como Huangascar.

  \item \textbf{Dashboard accionable:} Las recomendaciones de riego
    expresadas en horas concretas (basadas en ET\textsubscript{0} y
    déficit hídrico) son directamente utilizables por técnicos agrícolas
    sin formación en análisis de datos.
\end{enumerate}

\subsection{Limitaciones}

\begin{itemize}
  \item Los datos IoT son sintéticos; la implementación real requeriría
    la instalación de sensores físicos con comunicación LoRa o NB-IoT.
  \item Los modelos ML fueron entrenados sobre datos simulados, por lo
    que las métricas reportadas son orientativas. La validación con datos
    reales de campo es necesaria antes del despliegue productivo.
  \item El sistema opera en modo batch-diario; para alertas de helada en
    tiempo real se requeriría integración con streams de datos de
    pronóstico (SENAMHI).
\end{itemize}

\subsection{Trabajo futuro}

\begin{itemize}
  \item Integración de datos SENAMHI en tiempo real vía API REST para
    pronóstico de heladas 72 horas.
  \item Despliegue de nodos IoT físicos con sensores capacitivos de
    humedad de suelo (Decagon 5TM) y datalogger LoRaWAN.
  \item Modelos de deep learning (LSTM) para pronóstico de series
    temporales climáticas a 7 días.
  \item Extensión del sistema a otras culturas andinas: quinua, maíz
    morado, oca.
  \item Interfaz móvil para agricultores con alertas push vía WhatsApp
    o SMS.
\end{itemize}

% ============================================================
\section{Conclusiones}
% ============================================================

\begin{enumerate}
  \item Se diseñó e implementó exitosamente un pipeline Big Data completo
    para el monitoreo agroclimático de papa nativa en Huangascar, Yauyos,
    integrando 72\,397 registros de tres fuentes heterogéneas en un Data
    Lake AWS S3.

  \item El pipeline ETL sobre Apache Spark 3.4 en EMR 6.15 procesó los
    datos en menos de 3 minutos sobre un clúster de costo mínimo,
    demostrando la escalabilidad del enfoque distribuido para el contexto
    agrícola andino.

  \item Los cuatro índices agroclimáticos (IHB, ITT, IEH, ET\textsubscript{0})
    permiten caracterizar los riesgos de helada, tizón tardío y estrés
    hídrico con mayor sensibilidad que indicadores genéricos,
    identificando 18, 47 y 31 días críticos por año respectivamente.

  \item El regresor GBT Regressor obtuvo en la ejecución real sobre EMR un
    R\textsuperscript{2} de \textbf{0.6764} y un RMSE de \textbf{0.97 t/ha}
    para la predicción de rendimiento de papa nativa, con un rendimiento
    promedio predicho de 9.2 t/ha para la campaña 2025. El clasificador
    RandomForest identificó 288 días con alto riesgo de tizón (7.9\,\%
    del período histórico), revelando la naturaleza episódica del patógeno
    en el clima andino de Huangascar.

  \item El dashboard Streamlit integra de forma cohesiva los resultados
    del pipeline mediante 11 componentes visuales interactivos, incluyendo
    semáforos de riesgo, recomendaciones de riego basadas en
    ET\textsubscript{0} y predicciones por variedad, constituyendo una
    herramienta accionable para técnicos agrícolas.

  \item El sistema AgroSmart Andino valida la hipótesis central del
    proyecto: es técnicamente factible y económicamente viable implementar
    un sistema Big Data de monitoreo agroclimático de alta precisión para
    la pequeña agricultura andina utilizando herramientas de código
    abierto y servicios cloud de bajo costo.
\end{enumerate}

% ============================================================
\section*{Referencias}
% ============================================================
\addcontentsline{toc}{section}{Referencias}

\begin{thebibliography}{9}

\bibitem{midagri2023}
MIDAGRI. (2023). \textit{Estadísticas Agrícolas: Producción de Papa por
Regiones 2000--2022}. Ministerio de Desarrollo Agrario y Riego del Perú.
\url{https://www.midagri.gob.pe/portal/estadisticas}

\bibitem{senamhi2022}
SENAMHI. (2022). \textit{Boletín Climático Mensual — Región Lima}.
Servicio Nacional de Meteorología e Hidrología del Perú.
\url{https://www.senamhi.gob.pe}

\bibitem{kamilaris2017}
Kamilaris, A., \& Prenafeta-Boldú, F. X. (2018). Deep learning in
agriculture: A survey. \textit{Computers and Electronics in Agriculture},
147, 70--90. \url{https://doi.org/10.1016/j.compag.2018.02.016}

\bibitem{hargreaves1985}
Hargreaves, G. H., \& Samani, Z. A. (1985). Reference crop
evapotranspiration from temperature. \textit{Applied Engineering in
Agriculture}, 1(2), 96--99. \url{https://doi.org/10.13031/2013.26773}

\bibitem{nasapower}
Stackhouse, P. W. (2023). \textit{NASA POWER Project: Prediction of
Worldwide Energy Resource}. NASA Langley Research Center.
\url{https://power.larc.nasa.gov}

\bibitem{databricks2023}
Chambers, B., \& Zaharia, M. (2018). \textit{Spark: The Definitive Guide}.
O'Reilly Media.

\bibitem{cassandraguide}
Carpenter, J., \& Hewitt, E. (2022). \textit{Cassandra: The Definitive
Guide} (3rd ed.). O'Reilly Media.

\bibitem{streamlit2024}
Treuille, A., Kelly, T., \& Toivonen, S. (2024). \textit{Streamlit: The
fastest way to build data apps}. \url{https://streamlit.io/docs}

\bibitem{agrotec2021}
INIA. (2021). \textit{Tecnologías para el manejo integrado del cultivo de
papa en la sierra del Perú}. Instituto Nacional de Innovación Agraria.

\end{thebibliography}

% ── Anexos ──────────────────────────────────────────────────
\appendix
\section{Estructura de archivos del proyecto}

\begin{lstlisting}[style=bashstyle, caption={Árbol de archivos del proyecto}]
agrosmart_andino/
+-- README.md
+-- requirements.txt
+-- .env.example
+-- .gitignore
+-- data/
|   +-- raw/
|   |   +-- nasa_power_yauyos.csv      # 18,265 filas  1.6 MB
|   |   +-- sensores_iot.csv           # 52,632 filas  4.6 MB
|   |   +-- produccion_papa_yauyos.csv #  1,500 filas  0.1 MB
|   +-- gold/
|       +-- lecturas_climaticas.csv    #  3,653 filas
|       +-- alertas_activas.csv        #    500 filas
|       +-- recomendaciones_riego.csv  #    300 filas
|       +-- predicciones_rendimiento.csv #   6 filas
+-- scripts/
|   +-- data_acquisition/
|   |   +-- 01_download_nasa_power.py
|   |   +-- 02_generate_iot_data.py
|   |   +-- 03_prepare_midagri_data.py
|   |   +-- 04_validate_datasets.py
|   +-- aws/
|   |   +-- 05_upload_to_s3.py
|   |   +-- 06_create_emr_cluster.py
|   +-- cassandra/
|       +-- 07_create_schema.cql
|       +-- 08_load_data.py
|       +-- 09_setup_cassandra_cloud9.sh
+-- notebooks/
|   +-- agrosmart_pipeline.ipynb      # 38 celdas PySpark
+-- dashboard/
|   +-- app.py                        # Dashboard Streamlit
+-- docs/
    +-- informe_tecnico.tex           # Este documento
\end{lstlisting}

\section{Instrucciones de ejecución completa}

\begin{lstlisting}[style=bashstyle, caption={Flujo de ejecución paso a paso}]
# 1. Instalacion de dependencias
pip install -r requirements.txt

# 2. Adquisicion de datos (local)
python scripts/data_acquisition/01_download_nasa_power.py
python scripts/data_acquisition/02_generate_iot_data.py
python scripts/data_acquisition/03_prepare_midagri_data.py
python scripts/data_acquisition/04_validate_datasets.py

# 3. Carga a AWS S3
# Configurar credenciales en .env
python scripts/aws/05_upload_to_s3.py
python scripts/aws/06_create_emr_cluster.py

# 4. Pipeline PySpark (JupyterHub en EMR)
# Abrir notebooks/agrosmart_pipeline.ipynb en JupyterHub
# Ejecutar todas las celdas (Run All)

# 5. Descargar resultados Gold
aws s3 sync s3://agrosmart-andino-demo/gold/csv/ data/gold/

# 6. Configurar Cassandra en Cloud9
bash scripts/cassandra/09_setup_cassandra_cloud9.sh
python scripts/cassandra/08_load_data.py --host 127.0.0.1

# 7. Iniciar Dashboard
streamlit run dashboard/app.py
# Acceso: https://agrosmart-andino.streamlit.app/
\end{lstlisting}

\section{Configuración de credenciales AWS}

\begin{lstlisting}[style=bashstyle, caption={Archivo .env para AWS Academy}]
# Obtener de: AWS Academy > Learner Lab > AWS Details
AWS_ACCESS_KEY_ID=ASIA...
AWS_SECRET_ACCESS_KEY=...
AWS_SESSION_TOKEN=...
AWS_DEFAULT_REGION=us-east-1
S3_BUCKET=agrosmart-andino-demo
\end{lstlisting}

\end{document}
